% GENERATED FILE. Edit canonical lessons, then run npm run generate:books.
% canonical-source-hash: fnv1a64:4103372453b77a67
% canonical-lessons: GE-C523-fernbedienung, GE-C523-zappen, GE-C523-serie, GE-C523-staffel, GE-C523-einschaltquote

\chapter{Vor dem Fernseher}
\label{ch:ge-vor-dem-fernseher}
\hlchaptermodality{523}

\begin{chapteropening}
\textbf{By the end of this chapter:} \emph{I can talk about watching television.}

Complete the last lesson of chapter 523: I can talk about watching television.
\end{chapteropening}

\section[die Fernbedienung]{die Fernbedienung — remote control}

\label{lesson:GE-C523-fernbedienung}



\begin{quote}
Before the new one: say the German for publisher, then the German for translation.
\end{quote}

\subsection*{You'll want to know: die Fernbedienung}
\textbf{die Fernbedienung} — \textquotedblleft{}remote control\textquotedblright{}.

\begin{quote}
Die Fernbedienung liegt unter dem Sofa.
\end{quote}

\begin{quote}
\emph{The remote control is under the sofa.}
\end{quote}

Plural: \emph{die Fernbedienungen}.

\begin{cousinweb}
\emph{Fernbedienung} is built on \emph{fern} \textquotedblleft{}far\textquotedblright{}, from Old High German \emph{ferrana} (EtymDB), and \emph{Bedienung} \textquotedblleft{}operation, service\textquotedblright{}, from \emph{bedienen}.
\end{cousinweb}

\begin{grammarlens}[title={What you've built}]
One more word for telling, explaining and describing at B1.
\end{grammarlens}

\subsection*{Your turn}
\begin{itemize}
\raggedright
  \item \emph{Say it:} \emph{die Fernbedienung}
  \item \emph{Say it:} \emph{die Fernbedienung}, once more
  \item \emph{Say it:} a sentence of your own with \emph{die Fernbedienung}
  \item \emph{Recall:} say the German for publisher, then the German for translation, then say \emph{die Fernbedienung} again
\end{itemize}

\begin{tcolorbox}[breakable,colback=teal!4,colframe=teal!35!black,title={Before you move on}]
What does die Fernbedienung mean? (\textquotedblleft{}Remote control\textquotedblright{}.) Say it once more.
\end{tcolorbox}
\section[zappen]{zappen — to flick through channels}

\label{lesson:GE-C523-zappen}



\begin{quote}
Before the new one: say the German for translation, then the German for remote control.
\end{quote}

\subsection*{You'll want to know: zappen}
\textbf{zappen} — \textquotedblleft{}to flick through channels\textquotedblright{}.

\begin{quote}
Am Abend zappt er stundenlang durch die Programme.
\end{quote}

\begin{quote}
\emph{In the evening he flicks through the channels for hours.}
\end{quote}

Informal. Regular: \emph{hat gezappt}.

\begin{cousinweb}
\emph{zappen} was borrowed from English \emph{to zap}; EtymDB has no entry for the German verb.
\end{cousinweb}

\begin{grammarlens}[title={What you've built}]
One more word for telling, explaining and describing at B1.
\end{grammarlens}

\subsection*{Your turn}
\begin{itemize}
\raggedright
  \item \emph{Say it:} \emph{zappen}
  \item \emph{Say it:} \emph{zappen}, once more
  \item \emph{Say it:} a sentence of your own with \emph{zappen}
  \item \emph{Recall:} say the German for translation, then the German for remote control, then say \emph{zappen} again
\end{itemize}

\begin{tcolorbox}[breakable,colback=teal!4,colframe=teal!35!black,title={Before you move on}]
What does zappen mean? (\textquotedblleft{}To flick through channels\textquotedblright{}.) Say it once more.
\end{tcolorbox}
\section[die Serie]{die Serie — series, TV series}

\label{lesson:GE-C523-serie}



\begin{quote}
Before the new one: say the German for remote control, then the German for to flick through channels.
\end{quote}

\subsection*{You'll want to know: die Serie}
\textbf{die Serie} — \textquotedblleft{}series, TV series\textquotedblright{}.

\begin{quote}
Meine Schwester schaut gerade eine neue Serie über Ärzte.
\end{quote}

\begin{quote}
\emph{My sister is watching a new series about doctors.}
\end{quote}

Plural: \emph{die Serien}.

\begin{cousinweb}
\emph{Serie} comes from Latin \emph{seriēs} \textquotedblleft{}a row, a chain\textquotedblright{} (EtymDB), built on \emph{serere} \textquotedblleft{}to join, to bind\textquotedblright{}.
\end{cousinweb}

\begin{grammarlens}[title={What you've built}]
One more word for telling, explaining and describing at B1.
\end{grammarlens}

\subsection*{Your turn}
\begin{itemize}
\raggedright
  \item \emph{Say it:} \emph{die Serie}
  \item \emph{Say it:} \emph{die Serie}, once more
  \item \emph{Say it:} a sentence of your own with \emph{die Serie}
  \item \emph{Recall:} say the German for remote control, then the German for to flick through channels, then say \emph{die Serie} again
\end{itemize}

\begin{tcolorbox}[breakable,colback=teal!4,colframe=teal!35!black,title={Before you move on}]
What does die Serie mean? (\textquotedblleft{}Series, TV series\textquotedblright{}.) Say it once more.
\end{tcolorbox}
\section[die Staffel]{die Staffel — season (of a series); relay team}

\label{lesson:GE-C523-staffel}



\begin{quote}
Before the new one: say the German for to flick through channels, then the German for series.
\end{quote}

\subsection*{You'll want to know: die Staffel}
\textbf{die Staffel} — \textquotedblleft{}season (of a series); relay team\textquotedblright{}.

\begin{quote}
Die zweite Staffel ist noch spannender als die erste.
\end{quote}

\begin{quote}
\emph{The second season is even more exciting than the first.}
\end{quote}

Plural: \emph{die Staffeln}. In sport, a relay team.

\begin{cousinweb}
\emph{Staffel} first meant \textquotedblleft{}step, rung\textquotedblright{}; EtymDB traces it to Proto-Germanic \emph{stapulaz} \textquotedblleft{}post\textquotedblright{}. The sense of a season of episodes is recent.
\end{cousinweb}

\begin{grammarlens}[title={What you've built}]
One more word for telling, explaining and describing at B1.
\end{grammarlens}

\subsection*{Your turn}
\begin{itemize}
\raggedright
  \item \emph{Say it:} \emph{die Staffel}
  \item \emph{Say it:} \emph{die Staffel}, once more
  \item \emph{Say it:} a sentence of your own with \emph{die Staffel}
  \item \emph{Recall:} say the German for to flick through channels, then the German for series, then say \emph{die Staffel} again
\end{itemize}

\begin{tcolorbox}[breakable,colback=teal!4,colframe=teal!35!black,title={Before you move on}]
What does die Staffel mean? (\textquotedblleft{}Season (of a series); relay team\textquotedblright{}.) Say it once more.
\end{tcolorbox}
\section[die Einschaltquote]{die Einschaltquote — ratings, viewing figures}

\label{lesson:GE-C523-einschaltquote}



\begin{quote}
Before the new one: say the German for series, then the German for season.
\end{quote}

\subsection*{You'll want to know: die Einschaltquote}
\textbf{die Einschaltquote} — \textquotedblleft{}ratings, viewing figures\textquotedblright{}.

\begin{quote}
Die Show hatte eine hohe Einschaltquote.
\end{quote}

\begin{quote}
\emph{The show had high ratings.}
\end{quote}

Plural: \emph{die Einschaltquoten}. Built on \emph{einschalten} — to switch on.

\begin{cousinweb}
\emph{Einschaltquote} is built on \emph{einschalten} \textquotedblleft{}to switch on\textquotedblright{} and \emph{Quote} \textquotedblleft{}share\textquotedblright{}, which goes back to Latin \emph{quotus} \textquotedblleft{}how many\textquotedblright{} (EtymDB).
\end{cousinweb}

\begin{grammarlens}[title={What you've built}]
One more word for telling, explaining and describing at B1.
\end{grammarlens}

\subsection*{Your turn}
\begin{itemize}
\raggedright
  \item \emph{Say it:} \emph{die Einschaltquote}
  \item \emph{Say it:} \emph{die Einschaltquote}, once more
  \item \emph{Say it:} a sentence of your own with \emph{die Einschaltquote}
  \item \emph{Recall:} say the German for series, then the German for season, then say \emph{die Einschaltquote} again
\end{itemize}

\begin{tcolorbox}[breakable,colback=teal!4,colframe=teal!35!black,title={Before you move on}]
What does die Einschaltquote mean? (\textquotedblleft{}Ratings, viewing figures\textquotedblright{}.) Say it once more.
\end{tcolorbox}
